\documentclass[11pt]{article}

\usepackage{amsmath,amssymb}
\usepackage{booktabs}
\usepackage{geometry}
\usepackage{enumitem}
\usepackage{hyperref}
\usepackage{fancyhdr}
\usepackage{xcolor}

\geometry{margin=1in}
\hypersetup{colorlinks=true,linkcolor=blue,urlcolor=blue}
\setlist[itemize]{leftmargin=1.3em,itemsep=0.3em,topsep=0.25em}
\setlist[enumerate]{leftmargin=1.5em,itemsep=0.35em,topsep=0.25em}
\pagestyle{fancy}
\fancyhf{}
\fancyhead[L]{\small Expanded guide --- Part I, Part II, Fisher pathways}
\fancyhead[R]{\small Helena}
\fancyfoot[C]{\thepage}
\renewcommand{\headrulewidth}{0.4pt}

\title{Expanded guide to the Lyman-$\alpha$ $P_{1\mathrm{D}}$ calculator\\[0.35em]
\large Background, Part~I, Part~II, code specifications,\\
and pathways toward Fisher forecasts and derivatives}
\author{Helena}
\date{12 November 2026}

\begin{document}
\maketitle

\begin{abstract}
\noindent
This document expands the earlier outline into a full study guide.
It explains the scientific background of the Lyman-$\alpha$ forest power spectrum, why the calculator was built the way it was, what Part~I and Part~II contain in detail (equations, code paths, figures, and limits), and how one can move from this forward model toward a Fisher forecast, including concrete options for taking derivatives (finite differences, autodiff with JAX, and which model to differentiate first).
\end{abstract}

\paragraph{Repository.}
\texttt{fisher-forecasts-future-spectroscopic-survey}.
\textbf{NumPy package:} \texttt{p1d/}.
\textbf{JAX analytic path (phase~1):} \texttt{p1d\_jax/}.
\textbf{Technical notes:} \texttt{docs/part1.tex}, \texttt{docs/part2.tex}.

\paragraph{Fiducial point.}
$\Omega_m=0.31$, $\Omega_b=0.049$, $h=0.67$,
$n_s=0.965$, $A_s=2.1\times 10^{-9}$,
default $z=3$, $b_F=-0.14$, $\beta_F=1.4$.

\tableofcontents
\newpage

%======================================================================
\section{Scientific background}
%======================================================================

\subsection{What the Lyman-$\alpha$ forest is}

Light from a distant quasar travels through the intervening intergalactic medium (IGM).
Neutral hydrogen absorbs at the Lyman-$\alpha$ wavelength in the rest frame of the gas.
Because of cosmic expansion, absorption at different distances appears at different observed wavelengths, so a single quasar spectrum becomes a one-dimensional map of transmission along the line of sight.
That fluctuating absorption pattern is the Lyman-$\alpha$ forest.

The transmitted flux fraction is written $F=e^{-\tau}$, where $\tau$ is the optical depth.
On large scales the forest traces the underlying dark-matter density field; on small scales gas pressure, thermal broadening, and nonlinear mapping from density to flux matter.
Surveys such as eBOSS and DESI measure the clustering of this forest and use it for baryon acoustic oscillations and, potentially, for broadband cosmology and neutrino-mass constraints.

\subsection{What $P_{1\mathrm{D}}$ is and why we use it}

A single spectrum only resolves structure along the line of sight.
The natural auto-statistic of one spectrum is therefore the \textbf{one-dimensional flux power spectrum} $P_{1\mathrm{D}}(k_\parallel,z)$: how much power the flux contrast has as a function of the line-of-sight wavenumber.
In observations, $k_\parallel$ is quoted in velocity units ($\mathrm{s\,km}^{-1}$) and $P_{1\mathrm{D}}$ in $\mathrm{km\,s}^{-1}$.

This is not the same as the three-dimensional flux power $P_F(k,\mu)$, which needs many neighbouring sightlines or a volume so that transverse modes are defined.
Part~I focuses on $P_{1\mathrm{D}}$ because it is the simplest auto-spectrum of the forest and a natural first forward model for a forecast.
Part~II adds a three-dimensional measurement of $P_F(k,\mu)$ so that redshift-space anisotropy can be tested.

\subsection{Why a Fisher forecast needs a forward model}

A Fisher matrix forecast answers: given a survey and a model, how well can parameters be constrained?
It needs:
\begin{enumerate}
\item a \textbf{mean prediction} $\mathbf{m}(\boldsymbol{\theta})$ --- here a spectrum as a function of parameters $\boldsymbol{\theta}$;
\item a \textbf{covariance} $C$ of the data vector;
\item \textbf{derivatives} $\partial\mathbf{m}/\partial\theta_i$.
\end{enumerate}
The Fisher matrix is schematically
\begin{equation}
F_{ij}
=
\left(\frac{\partial\mathbf{m}}{\partial\theta_i}\right)^{\!T}
C^{-1}
\left(\frac{\partial\mathbf{m}}{\partial\theta_j}\right)
\end{equation}
(plus terms if the covariance also depends on parameters).
Parts~I and~II only build $\mathbf{m}(\boldsymbol{\theta})$ for the flux power.
They do not yet build $C$, survey windows, noise, or $F_{ij}$.
That is intentional: an undefined flux model would make later Fisher work ambiguous.

\subsection{Why two flux paths?}

There is no unique way to turn linear $P_\mathrm{CDM}$ into forest flux power.
\begin{itemize}
\item On the largest scales, a linear bias plus coherent peculiar velocities (Kaiser / McDonald) is standard and differentiable.
\item On smaller scales the map $F=e^{-\tau(\Delta)}$ is nonlinear, so a purely linear bias cannot be the whole story.
\end{itemize}
Part~I therefore implements \textbf{both} an analytic Kaiser path and an optical-depth (FGPA) path, fed by the \textbf{same} $P_\mathrm{CDM}$.
When they disagree, the disagreement is in the flux model, not in a hidden difference of the matter spectrum.
That design choice is the backbone of the whole calculator.

%======================================================================
\section{Glossary}
%======================================================================

\begin{description}[style=nextline,leftmargin=1.2em]
\item[$P_\mathrm{CDM}$]
Linear cold-dark-matter matter power spectrum $P(k,z)$.
Input to everything. Not the flux power.

\item[$P_{1\mathrm{D}}$]
One-dimensional power along a line of sight, function of $k_\parallel$ only.

\item[$P_F(k,\mu)$]
Three-dimensional flux power; $\mu=k_\parallel/k$.

\item[Fiducial]
Baseline cosmology listed on the title page; ``vary $A_s$'' means relative to this point.

\item[Kaiser]
$P_F=b_F^{2}(1+\beta_F\mu^{2})^{2}P_\mathrm{CDM}$ with fixed $b_F=-0.14$, $\beta_F=1.4$.

\item[FGPA]
Fluctuating Gunn--Peterson Approximation: $\tau=A\Delta^{\beta}$, $F=e^{-\tau}$.

\item[Lognormal]
$\Delta=\exp(\delta_G-\tfrac12\mathrm{Var}(\delta_G))$; positive density from a Gaussian contrast.

\item[Jeans / pressure]
$\exp(-2 x_J^{2} k^{2})$ multiplying $P_\mathrm{CDM}$; $x_J=100$ or $50\,\mathrm{kpc}/h$.

\item[Arinyo]
Arinyo-i-Prats et al.\ (2015): hydro-calibrated fitting formula for nonlinear $P_F(k,\mu)$.

\item[Fractional change]
$(P_{\mathrm{new}}-P_{\mathrm{fid}})/P_{\mathrm{fid}}$. Not a fit.
\end{description}

\newpage
%======================================================================
\section{PART I --- baseline calculator}
%======================================================================

\subsection{Goal and scope}

\textbf{Goal.} Map $(n_s,A_s,P_\mathrm{CDM})\to P_{1\mathrm{D}}(k_\parallel,z)$ at one redshift.

\textbf{In scope.} Eisenstein--Hu (or table) matter spectrum; Kaiser projection; FGPA sightlines; check plots 01--08.

\textbf{Out of scope.} Cross-power, covariance, survey window, noise, Fisher matrix.

\subsection{Units and cosmology}

The conversion between comoving and forest velocity units is
\begin{equation}
\frac{\mathrm{d}v}{\mathrm{d}\chi}=\frac{H(z)}{(1+z)\,h}
\quad[\mathrm{km\,s}^{-1}\text{ per }\mathrm{Mpc}/h].
\end{equation}
At the fiducial $z=3$ this is $\approx 113$.
Then $k\,[h\,\mathrm{Mpc}^{-1}]=k\,[\mathrm{s\,km}^{-1}]\,\mathrm{d}v/\mathrm{d}\chi$ and
$P_{1\mathrm{D}}\,[\mathrm{km\,s}^{-1}]=P_{1\mathrm{D}}\,[\mathrm{Mpc}/h]\,\mathrm{d}v/\mathrm{d}\chi$.
The window $10^{-3}$--$0.02\,\mathrm{s\,km}^{-1}$ is about $0.11$--$2.3\,h\,\mathrm{Mpc}^{-1}$.

Why fix $\Omega_m$, $\Omega_b$, $h$?
So that variations of $n_s$ and $A_s$ can be read cleanly in the check plots.
Background expansion still enters through $H(z)$ and growth $D(z)$.

\subsection{Building $P_\mathrm{CDM}$ --- why Eisenstein--Hu}

A full Boltzmann code (CAMB/CLASS) is the gold standard, but a Fisher loop may evaluate the forward model thousands of times.
Eisenstein \& Hu (1998) give a closed fitting formula for the transfer function with baryon acoustic oscillations.
It is accurate enough for this stage (median ratio to CAMB $\sim 1.01$ on $10^{-3}$--$10\,h\,\mathrm{Mpc}^{-1}$ at $z=0$ in our check) and makes $n_s$ and $A_s$ trivial to vary.

Primordial amplitude $A_s$ is defined at pivot $k_\mathrm{pivot}=0.05\,\mathrm{Mpc}^{-1}$.
Changing $n_s$ tilts the spectrum about that pivot; changing $A_s$ rescales the whole spectrum.

\textbf{Code.}
\texttt{p1d/eisenstein\_hu.py}:
\texttt{transfer}, \texttt{growth}, \texttt{hubble\_kms}, \texttt{linear\_matter\_power}.
Growth is normalised so $D(a)\to a$ in matter domination, with a numerical integral on a log-$a$ grid.

\textbf{Alternatives implemented for flexibility.}
\begin{itemize}
\item \texttt{TabulatedMatterPower} (\texttt{matter\_table.py}): load an external linear table; rescale in $A_s$, $n_s$, and growth without rebuilding EH.
Used so a CLASS/CAMB file can replace EH later without rewriting the flux layer.
\item $\alpha\beta\gamma\delta$ deformation (\texttt{abgd.py}): multiply EH by
$\exp(\alpha L+\beta L^{2}+\gamma L^{3}+\delta L^{4})$, $L=\ln(k/k_p)$, $k_p=1\,h\,\mathrm{Mpc}^{-1}$.
At the pivot the factor is 1.
\textbf{This warp is an internal shape knob, not taken from a Ly$\alpha$ paper}; EH remains the physical base.
\end{itemize}

\subsection{Method 1 --- Kaiser path (why and how)}

\subsubsection{Physics}

On scales where the density contrast is small and velocities are coherent, the forest can be treated as a linearly biased tracer.
Kaiser (1987) derived the redshift-space factor for galaxies; McDonald (2003) wrote the analogous large-scale forest model.
The three-dimensional flux power is
\begin{equation}
P_F(k,\mu)=b_F^{2}\,(1+\beta_F\mu^{2})^{2}\,P_\mathrm{CDM}(k),
\qquad\mu=k_\parallel/k.
\end{equation}
We fix $b_F=-0.14$ and $\beta_F=1.4$ at every redshift.
These are a round form of early DESI large-scale measurements (Gordon et al.\ 2023).
The code does \emph{not} compute $\beta=f/b$ from first principles, and it does \emph{not} apply the $(1+z)^{2.9}$ bias evolution used in eBOSS/DESI analyses.
McDonald's extra scale-dependent corrections (nonlinear growth, pressure, nonlinear velocities) are omitted: this path is linear Kaiser only.

Why negative $b_F$?
An overdensity raises the neutral fraction and optical depth and therefore \emph{lowers} transmission, so the flux contrast anti-correlates with density on large scales.

\subsubsection{Projection to $P_{1\mathrm{D}}$}

A spectrum measures only $k_\parallel$, so
\begin{equation}
P_{1\mathrm{D}}(k_\parallel)
=
\int_{k_\parallel}^{k_\mathrm{max}}
\frac{k\,\mathrm{d}k}{2\pi}\,
P_F\!\left(k,\mu=\tfrac{k_\parallel}{k}\right),
\end{equation}
with $k_\mathrm{max}=30\,h\,\mathrm{Mpc}^{-1}$ and a trapezoidal rule on $\sim 800$ log-$k$ nodes.
The lower limit is kinematic: a mode of wavenumber $k$ cannot contribute to a larger $k_\parallel$ than $k$.
The upper limit is past the plotted forest window because high-$k$ power still contributes to lower $k_\parallel$.
There is no random field: the result is deterministic and smooth in $n_s$ and $A_s$, which is ideal for derivatives.

\textbf{Code.}
\texttt{p1d/flux.py}:
\texttt{flux\_power\_3d}, \texttt{flux\_power\_1d}, \texttt{p1d\_skm},
plus unit helpers \texttt{k\_skm\_to\_hmpc}, \texttt{p1d\_h\_to\_kms}.
Optional argument \texttt{matter\_power=} replaces the internal EH call --- this is the hook Part~II uses for Jeans-smoothed tables.

\textbf{Entry point for forecasts (analytic).}
\texttt{p1d\_skm(k\_skm, z, ns, As, omega\_m, omega\_b, h, ...)}.

\subsection{Method 2 --- FGPA sightlines (why and how)}

\subsubsection{Why a second path?}

Kaiser assumes linear bias.
The actual forest observable is transmission $F=e^{-\tau}$ with $\tau$ a nonlinear function of density.
The sightline path builds that map explicitly so we can see where Kaiser is only an amplitude (or fails).

\subsubsection{Step-by-step algorithm}

\begin{enumerate}
\item \textbf{Isotropic 1D matter projection}
(\texttt{projection.project\_p1d}):
\[
P_{1\mathrm{D,matter}}(k_\parallel)
=
\frac{1}{2\pi}\int_{k_\parallel}^{q_\mathrm{max}} q\,P_\mathrm{CDM}(q)\,\mathrm{d}q.
\]
Still no flux bias --- only the geometry of projecting an isotropic 3D field onto one line.

\item \textbf{Gaussian realisation}
(\texttt{skewers.generate\_gaussian\_skewers}):
draw complex Fourier modes with variance fixed by that 1D matter power; zero mode set to zero; inverse real FFT $\to$ $\delta_G$.
Distinct lines are independent.
A single line has the same 1D power as a row through an isotropic Gaussian volume, but \emph{no} transverse correlations --- so $P_F(k,\mu)$ cannot be measured from these lines alone.
That is why Part~II needed a cube.

\item \textbf{Lognormal density}
(\texttt{fgpa\_field.lognormal\_density}):
\[
\Delta=\exp\!\left(\delta_G-\tfrac12\mathrm{Var}(\delta_G)\right).
\]
A Gaussian contrast can satisfy $\delta_G<-1$, so $1+\delta_G$ can be negative.
Gas density cannot be negative, and $\Delta^{\beta}$ is undefined for $\Delta\le 0$.
The exponential forces positivity; subtracting half the variance keeps $\langle\Delta\rangle$ near one.
References: Coles \& Jones (1991); Bi \& Davidsen (1997).

\item \textbf{Optical depth (FGPA)}
(\texttt{prepare\_tau\_unit}, \texttt{calibrate\_mean\_flux}):
\[
\beta=2-0.7(\gamma-1),\qquad
\tau=A\Delta^{\beta},\qquad
F=e^{-\tau}.
\]
Default $\gamma=1.5$ $\Rightarrow$ $\beta=1.65$ (recombination physics; \textbf{not} Kaiser $\beta_F$).
Solve $A$ so $\langle F\rangle$ equals a target (default $0.80$) using a Brent root find on $\ln A$.
For the Part~I comparison ensemble ($256$ lines, $L=200\,\mathrm{Mpc}/h$, $N=512$), $A\approx 0.485$.

\item \textbf{Flux power}
$\delta_F=F/\langle F\rangle-1$; periodogram averaged over lines; convert to $\mathrm{km\,s}^{-1}$.
\end{enumerate}

Optional pixel Gaussians for pressure and thermal broadening exist (\texttt{smooth\_periodic}) but are left at zero width in Part~I and Part~II figures.
A LOS displacement can be passed for peculiar velocities, but Part~I comparison figures are \textbf{real space}.

\textbf{Code entry point.}
\texttt{run\_fgpa\_skewers(TabulatedMatterPower, SkewerConfig)} in \texttt{fgpa\_field.py}.

\subsubsection{Why absolute amplitudes disagree with Kaiser}

Linearising FGPA for small contrast gives $\delta_F\approx -A\beta\,\delta$.
On the comparison lines $A\beta\approx 0.80$, several times $|b_F|=0.14$.
So the two paths are normalised by different conditions: a large-scale bias (+ RSD) versus a prescribed mean transmission in real space.
Also, without pressure smoothing the lognormal variance on those lines is $\sim 19$ --- far outside the linear regime --- so one should not expect absolute agreement.

\subsection{Part I figures in detail}

\begin{center}
\small
\begin{tabular}{@{}cp{0.82\textwidth}@{}}
\toprule
Fig & Narrative \\
\midrule
01 & $P_\mathrm{CDM}(k)$ for different $n_s$ and $A_s$: tilt vs overall amplitude. \\
02 & Same variations after Kaiser projection to $P_{1\mathrm{D}}$ (absolute units). \\
03 & $(P-P_{\mathrm{fid}})/P_{\mathrm{fid}}$:
$+1\%$ in $A_s$ $\to$ flat $+1\%$;
$\Delta n_s=+0.01$ $\to$ scale-dependent $\sim 2$--$4\%$.
Shows derivatives are well behaved on the Kaiser path. \\
04 & Redshift dependence at fixed cosmology (growth + unit conversion). \\
05 & $\alpha\beta\gamma\delta$ shape deformation; fractional change vanishes at $k_p=1\,h\,\mathrm{Mpc}^{-1}$. \\
06 & Table rescaling agrees with direct EH at $\sim 10^{-6}$. \\
07 & Full pipeline: matter, 1D matter, Kaiser vs FGPA absolute $P_{1\mathrm{D}}$, one skewer.
Absolute curves do not overlie. \\
08 & Divide each flux spectrum by its value at $0.005\,\mathrm{s\,km}^{-1}$.
\textbf{No fit} --- only renormalisation.
Shapes agree; wiggles = 256-line noise. \\
\bottomrule
\end{tabular}
\end{center}

\textbf{Reproduce:} \texttt{python make\_checks.py}

\subsection{Part I takeaway}

Two forward models from one $P_\mathrm{CDM}$.
Kaiser is the natural differentiable forecast engine.
FGPA sightlines are the physics check.
Shapes agree when amplitude is removed; absolute offset is expected.

\newpage
%======================================================================
\section{PART II --- six October extensions}
%======================================================================

\subsection{Implementation rule (why nothing was rewritten)}

The October requests were extensions, not a replacement of Part~I.
Keeping \texttt{p1d\_skm} and \texttt{run\_fgpa\_skewers} bit-compatible with October defaults means:
\begin{itemize}
\item regressions are easy (\texttt{pytest});
\item new physics is opt-in (pass a Jeans table, or call \texttt{box3d});
\item the Fisher path can stay on the analytic function without inheriting sightline noise.
\end{itemize}

\subsection{1. Pressure / Jeans filtering}

\textbf{Ask:} Gaussian pressure filtering on the gas.

\textbf{Why this form:} Arya et al.\ (2022) use $\exp(-2 x_J^{2} k^{2})$ on the linear spectrum in their lognormal forest model.
We multiply the \emph{stored} $P_\mathrm{CDM}$ before either flux path, rather than only smoothing pixels after the fact, so Kaiser and FGPA see the same cutoff.

\textbf{Code.}
\texttt{p1d/jeans.py}:
\texttt{jeans\_suppression(k, x\_j)} and \texttt{with\_jeans(table, x\_j)} (new table; input untouched; copies $A_s$, $n_s$, pivot).
Lengths: $100$ and $50\,\mathrm{kpc}/h$.

\textbf{Results (figs 09, 11).}
Large scales unchanged; $100\,\mathrm{kpc}/h$ cuts first.
Sightline flux falls less steeply than Kaiser at high $k$ because the lognormal regenerates small-scale power.

\subsection{2. Lognormal for weakly nonlinear dark matter}

\textbf{Ask:} use lognormal to capture weakly nonlinear DM.

\textbf{Clarification that mattered in the meeting:}
lognormal was \textbf{already} the density for $\tau$ in Part~I.
Part~II did not invent a new density model; it \textbf{diagnosed} Gaussian vs lognormal on short, highly resolved lines ($L=40\,\mathrm{Mpc}/h$, $N=1024$) so the Nyquist mode sits past the Jeans cutoff.

\textbf{Result (fig 10).}
Without Jeans, $\sim 1/3$ of Gaussian pixels have $1+\delta\le 0$ (unusable for FGPA).
At $100\,\mathrm{kpc}/h$, $\mathrm{Var}(\Delta)\sim 2$: weakly nonlinear.
Lognormal power exceeds Gaussian power at high $k$.

\subsection{3. More spectra / noise}

\textbf{Ask:} more sightlines; check that noise diminishes.

\textbf{Result (fig 12).}
From 256 to 10\,000 lines, scatter falls roughly as $N^{-1/2}$.
The Kaiser offset remains at 10\,000 lines $\Rightarrow$ model normalisation, not sample variance.

\subsection{4. Small scales at $z\sim 5$}

\textbf{Ask:} spectra to $k\sim 0.2\,\mathrm{s\,km}^{-1}$ at $z\sim 5$.

\textbf{Implementation choices.}
Short lines so Nyquist exceeds $0.2\,\mathrm{s\,km}^{-1}$; Kaiser integral to $80\,h\,\mathrm{Mpc}^{-1}$.
Mean transmission from Faucher-Gigu\`ere $\tau_\mathrm{eff}=0.0018(1+z)^{3.92}$ $\Rightarrow$ $\langle F\rangle\approx 0.13$.
\textbf{Extrapolated} past the fit range $z\le 4.2$.
Chosen instead of $0.80$ because a transparent forest is not realistic at $z=5$.

\textbf{Result (fig 13).}
Reach the requested scale; sightline power $\gg$ Kaiser because of the normalisation mismatch.

\subsection{5--6. 3D cube and Arinyo comparison}

\textbf{Ask:} simulate in 3D to measure $P_F(k,\mu)$; compare to Arinyo-i-Prats.

\textbf{Why 1D lines are not enough.}
Independent skewers lack transverse modes and cannot source LOS velocities from perpendicular modes.
Kaiser \emph{assumes} the $\mu$ factor; the cube \emph{measures} it.

\textbf{Implementation} (\texttt{p1d/box3d.py}, \texttt{p1d/arinyo.py}).
One cube $40\,\mathrm{Mpc}/h$, $128^{3}$, $z=3$, Jeans $100\,\mathrm{kpc}/h$.
Gaussian field with exact rms amplitudes + random phases.
Lognormal FGPA; linear continuity velocity; CIC shift of $\tau$; solve $A$ for $\langle F\rangle=0.696$ (Arinyo's $z=3$ value, intentionally different from Part~I's $0.80$).
Bin in $(k,\mu)$.
Estimator check recovers analytic Kaiser to $\sim 10^{-13}$.

\textbf{Four curves on fig 14 (each panel).}
\begin{enumerate}
\item Solid blue --- analytic Kaiser on the screened spectrum (formula, not a measurement).
\item Dashed green --- Arinyo published $z=3$ formula, not refitted.
\item Dotted orange --- cube flux in real space.
\item Solid red --- cube flux after linear RSD.
\end{enumerate}
Left: $\mu<0.25$ (transverse). Right: $\mu>0.75$ (parallel).

\textbf{Interpretation.}
Red above orange on the parallel panel and below on the transverse panel $\Rightarrow$ Kaiser sense of anisotropy from a measurement.
Shape tracks Arinyo; amplitude lies above it.
Limits: one small box; Nyquist $\sim 10\,h\,\mathrm{Mpc}^{-1}$; linear velocities only; no thermal broadening; Arinyo not refitted; lognormal $\neq$ hydro.

\textbf{Overleaf.}
Part~II note, section ``Three-dimensional flux power'', Figure~14.

\textbf{Reproduce Part~II.}
\begin{verbatim}
python make_extension_plots.py   # 09-12
python make_meeting_plots.py     # 13-14
python -m pytest -q
\end{verbatim}

\subsection{Code map}

\begin{center}
\small
\begin{tabular}{@{}p{0.42\textwidth}p{0.50\textwidth}@{}}
\toprule
Operation & Location \\
\midrule
EH $P_\mathrm{CDM}$, growth & \texttt{p1d/eisenstein\_hu.py} \\
Matter table & \texttt{p1d/matter\_table.py} \\
$\alpha\beta\gamma\delta$ & \texttt{p1d/abgd.py} \\
Kaiser $P_{1\mathrm{D}}$ & \texttt{p1d/flux.py} $\to$ \texttt{p1d\_skm} \\
1D projection / skewers & \texttt{projection.py}, \texttt{skewers.py} \\
FGPA sightlines & \texttt{fgpa\_field.py} $\to$ \texttt{run\_fgpa\_skewers} \\
Jeans & \texttt{jeans.py} $\to$ \texttt{with\_jeans} \\
Arinyo formula & \texttt{arinyo.py} \\
3D box & \texttt{box3d.py} \\
JAX Kaiser path (phase 1) & \texttt{p1d\_jax/} $\to$ \texttt{p1d\_skm} \\
\bottomrule
\end{tabular}
\end{center}

\newpage
%======================================================================
\section{Toward a Fisher forecast}
%======================================================================

\subsection{What is still missing}

\begin{enumerate}
\item Choose the \textbf{data vector}: e.g.\ $P_{1\mathrm{D}}(k_\parallel)$ in redshift bins; later $P_F(k,\mu)$ and/or cross-power $P_x$ with quasars/galaxies.
\item Build a \textbf{covariance} $C$ (Gaussian + sample variance + noise + resolution; possibly from the lognormal ensemble or an analytic model such as Arya et al.\ on covariances).
\item Include \textbf{survey realism}: area, tracer densities, spectrograph resolution, noise.
\item Compute \textbf{derivatives} $\partial\mathbf{m}/\partial\theta_i$.
\item Assemble $F_{ij}$ and report parameter errors / contours.
\end{enumerate}

\subsection{Which forward model should enter Fisher first?}

A practical order:
\begin{enumerate}
\item \textbf{Analytic Kaiser \texttt{p1d\_skm}} --- smooth, fast, already matches the structure of a Fisher mean model. Best first target for derivatives.
\item Optionally add IGM / bias parameters ($b_F$, $\beta_F$, $x_J$, mean-$F$ amplitude) once the cosmological derivatives work.
\item Use FGPA sightlines for \textbf{validation} of the mean model and for covariance estimation, not as the primary differentiable mean (unless you average enough lines that noise is negligible).
\item Use Arinyo or the cube only if the forecast data vector is 3D $P_F(k,\mu)$; both need more work (refit / larger volume) before they are forecast-ready.
\end{enumerate}

\subsection{How to take derivatives --- options}

\subsubsection{Finite differences (works today on NumPy \texttt{p1d})}

For each parameter $\theta$, evaluate
\begin{equation}
\frac{\partial\mathbf{m}}{\partial\theta}
\approx
\frac{\mathbf{m}(\theta+\epsilon)-\mathbf{m}(\theta-\epsilon)}{2\epsilon}.
\end{equation}
\textbf{Pros:} no code rewrite; works on the existing package.
\textbf{Cons:} step-size tuning; $2N_{\mathrm{param}}$ evaluations; noisy if $\mathbf{m}$ is stochastic (sightlines).
\textbf{Recommendation:} use only on \texttt{p1d\_skm} (deterministic). Figure~03 already shows that a $1\%$ $A_s$ shift produces a clean flat fractional response --- a good sanity check for $\partial/\partial A_s$.

\subsubsection{Autodiff with JAX (started as \texttt{p1d\_jax})}

Phase~1 ports Eisenstein--Hu and \texttt{p1d\_skm} to JAX so that
\begin{verbatim}
jax.grad / jax.jacfwd / jax.jacrev
\end{verbatim}
can differentiate the spectrum with respect to $n_s$, $A_s$, and later other parameters, with optional \texttt{jit}.
\textbf{Pros:} exact derivatives (to machine precision of the AD rules); one evaluation of the forward + VJP/JVP; natural fit to Fisher.
\textbf{Cons:} FGPA root-find for $A$, interpolators, and discrete FFTs need care; sightline randomness needs explicit keys.

\textbf{Suggested pathway.}
\begin{enumerate}
\item Finish validating \texttt{p1d\_jax.p1d\_skm} against NumPy (tests in \texttt{tests/test\_p1d\_jax.py}).
\item Differentiate wrt $(n_s,A_s)$ first; compare to finite differences.
\item Add Jeans length and $b_F$, $\beta_F$ as differentiable parameters.
\item Only later consider a differentiable FGPA mean-flux solve (replace Brent with a fixed-iteration Newton) if the forecast mean model needs it.
\end{enumerate}

\subsubsection{Emulator route}

Train a small emulator of $P_{1\mathrm{D}}(k\,|\,\boldsymbol{\theta})$ on a grid of Kaiser (or hydro) evaluations, then differentiate the emulator.
\textbf{Pros:} very fast in a Fisher/MCMC loop; can ingest expensive models.
\textbf{Cons:} training set design; interpolation error becomes part of the theory systematic.

\subsection{Parameter sets to discuss with collaborators}

\begin{itemize}
\item \textbf{Cosmology first:} $n_s$, $A_s$ (and later $\Omega_m$, $h$, $m_\nu$, \ldots).
\item \textbf{Bias / RSD:} $b_F$, $\beta_F$ --- free vs fixed vs Gaussian priors from DESI.
\item \textbf{IGM:} $x_J$, $\gamma$ or $\beta$, mean-$F(z)$ or $A$.
\item \textbf{Redshift:} independent $P_{1\mathrm{D}}$ bins vs a smooth $\langle F\rangle(z)$ model.
\end{itemize}

\subsection{Recommended near-term plan}

\begin{enumerate}
\item Lock the mean model for Fisher~v1: Kaiser \texttt{p1d\_skm} (NumPy or JAX), with a decision on whether $b_F$ is fixed.
\item Implement finite-difference derivatives and a toy diagonal covariance; produce a first Fisher ellipse for $(n_s,A_s)$ as a pipeline test.
\item Switch derivatives to \texttt{p1d\_jax} once tests match NumPy.
\item In parallel, decide $\langle F\rangle(z=5)$ and whether Jeans is fixed or a nuisance.
\item Only then expand to cross-power and a realistic covariance.
\end{enumerate}

%======================================================================
\section{Open decisions (ask collaborators)}
%======================================================================

\begin{enumerate}
\item Normalise Kaiser to FGPA/cube at large scales, or keep $b_F=-0.14$?
\item Which $\langle F\rangle$ at $z=5$?
\item Fix, vary, or drop explicit Jeans?
\item Refit Arinyo, keep published, or move to hydro/emulator?
\item Larger 3D box before trusting $P_F(k,\mu)$ in a forecast?
\item Thermal broadening / nonlinear RSD now or later?
\item Fisher data vector: $P_{1\mathrm{D}}$ only, or also $P_F(k,\mu)$ / $P_x$?
\item Derivatives: finite difference first, or JAX autodiff first?
\end{enumerate}

%======================================================================
\section{One-page recap}
%======================================================================

\textbf{Background.}
The forest is IGM absorption along quasar sightlines; $P_{1\mathrm{D}}$ is its 1D auto-spectrum and a natural first Fisher observable.

\textbf{Part~I.}
Same $P_\mathrm{CDM}$ $\to$ Kaiser $P_{1\mathrm{D}}$ (analytic) and FGPA sightlines (lognormal + $\tau=A\Delta^{\beta}$).
Shapes agree when amplitude is divided out; absolute offset is a normalisation difference.

\textbf{Part~II.}
Jeans on $P_\mathrm{CDM}$; Gaussian vs lognormal diagnostic; noise $\propto N^{-1/2}$ but Kaiser offset remains; $z=5$ reaches $0.2\,\mathrm{s\,km}^{-1}$; one cube measures $P_F(k,\mu)$ with linear RSD; shape vs Arinyo ok, amplitude high.

\textbf{Fisher next.}
Start from differentiable Kaiser \texttt{p1d\_skm}; build covariance and survey terms; take derivatives by finite difference and/or JAX (\texttt{p1d\_jax}); expand parameters and cross-power after the pipeline works end-to-end.

\end{document}
